\documentclass[11pt]{article}
\usepackage{xcolor} %textcolor
\usepackage{amsmath}
\begin{document}


\section{Method}
Simple methods that do not require explanation in \texttt{summary\_examples\_multichannel.mlx}
\begin{itemize}
	\item 0.2 ISI
	\item 
\end{itemize}
\subsection{ spike detection (0.2. in \texttt{summary\_examples\_multichannel.mlx})}
In order to detect neural spikes and build spike trains, we utilize a previous reported MATLAB spike-sorting algorithm~\cite{le_floch_stretchable_2022}.
The signal is first filtered bandpass filtered~(\textcolor{red}{300 -- 6000 Hz or 100 -- 6000 Hz}) using the MATLAB function \texttt{bandpass}.
Then any timepoint where the signal amplitude is lower than -\textcolor{red}{5}*standard deviation of the filtered  signal is considered as the peak of a spike.
We slice \textcolor{red}{2} milliseconds back and forth from the peak to make a spike waveform.

\subsection{burst detection}
For burst detection in spike trains, we follow~\cite{fair_electrophysiological_2020} to adopt MaxInterval method. 
It defines bursts using five pre-defined threshold parameters\textemdash
max ISI at start of burst,
max ISI in burst,
min burst duration,
min interburst interval,
and
min number of spikes in burst.
Any series of spikes that satisfies all of these thresholds is considered as a burst. We set the threshold as \textcolor{red}{300, 301, 10, 50, 200}, respectively.

\subsection{spike sorting through PCA and k-means~(0.4, 0.5 in \texttt{summary\_examples\_multichannel.mlx})}
In order to sort the detected spike waveforms that originates from different neurons, we apply dimension reduction and clustering algorithms.
We first use principal component analysis (MATLAB function \texttt{pca} to reduce each high-dimensional spike waveform into vectors residing in a \textcolor{red}{two}-dimensional feature space.
Then we cluster these reduced vectors using k-means clustering algorithm (MATLAB function \texttt{kmeans}), which clusters the vectors into the pre-specified number of groups.
We set the number of groups as \textcolor{red}{3}.

\subsection{Spike sorting quality metrics}
In order to assess how each spike clusteris well-separated from the spikes outside of the cluster, we adopt two measurements, namely L-ratio and isolation distance. 
Both measurements build on Malhalanobis distance between the $C$-th cluster and an $i$-th spike $x_i$ on the \textcolor{red}{two}-dimensional feature space.
The Mahalanobis distance accounts for the correlation between dimensions when measuring the distance, and is defined as follows:
\begin{align*}
	D_{i,C}:= (x_i-\mu_c)^\top
	\Sigma_c^{-1}
	(x_i-\mu_C)
\end{align*}
where $\mu_C$ and $\Sigma_C$ are the mean vector and the covariance matrix of the spike vectors in cluster $C$, respectively.
L-ratio assumes that the distribution of the spikes in the cluster $C$ follows a multivariate normal distribution, and is calculated as follows:
\begin{align*}
	L_{ratio} :=
	\frac{\sum_{i \notin C}
	\bigl(1-CDF_{\chi^2}(D_{i,C}^2)
	\bigr)}
	{\text{number of spikes in the cluster C}},n
\end{align*}where $i \notin C$ indicates the set of spikes which do not belong the the cluster $C$, and $CDF_{\chi^2}$ presents the cumulative distribution function of the chi-square  distribution with degree of freedom $=$ \textcolor{red}{2}.
\textcolor{blue}{jongmin: degree of freedom should equal to the number of principal components}



\subsection{3d network graph}


In order to measure the connectivity between electrodes, we follow the synchronized score approach of \cite{lam_tissue-specific_2019}.
To calculate the synchronized score between two spike trains, we first measure SPIKE-distance~\cite{kreuz_monitoring_2013} using Python environment~(version 3.9.2) and Python package \textsc{Pyspike}~(version 0.7.0)~\cite{mulansky_pyspikepython_2016}.
\textcolor{red}{The SPIKE-distance is normalized to average of 1000 randomly generated baseline values; each baseline value is the SPIKE-distance between two uniformly randomly generated spike trains.}
Then, synchronized scores are defined as SPIKE-distance values subtracted from 1.
These sync scores range from 0 to 1, where a higher score indicates a high degree of synchrony.




Community structure on the electrodes is examined by graph-theoretic Louvain algorithm~\cite{PhysRevE.69.026113, blondel_fast_2008} which is based on modularity optimization.
In the analysis, we model the electrodes as nodes~(e.g., circles in the network map), whose 3d positions represent the position of the actual electrode. 
We set synchronized scores between every pair of active electrodes as edge~(e.g., lines in the network map) weights between the corresponding nodes.
The edges with synchronized scores less than \textcolor{red}{0.5} were filtered out.
The result of the Louvain algorithm is visualized through color and size, using the Python package \texttt{plotly}~(version  5.9.0).
Nodes with the same color (i.e., electrodes) belong to the same community network.
Synchronization scores are color-mapped, ranged from \textcolor{red}{0 (blue) to 1 (red)}.
In addition, the larger the number of connected electrodes becomes, the larger the size of the node becomes.

\subsection{Figures from }
\bibliographystyle{apalike}
\bibliography{reference}
\end{document}
